\documentclass[11pt]{article}

\usepackage[margin=1in]{geometry}
\usepackage{amsmath,amssymb,amsthm}
\usepackage{graphicx}
\usepackage{hyperref}
\usepackage{verbatim}

\newtheorem{theorem}{Theorem}
\newtheorem{remark}{Remark}

\title{A Finite-Coherence Counterexample to a Global Pressure Claim}
\author{Automatic Math Discovery Smoke Test}
\date{June 9, 2026}

\begin{document}
\maketitle

\section*{Source statement}

This packet concerns arXiv:2509.13121, \emph{On the Fixed Point Property in Reflexive Banach Spaces}, by Faruk Alpay and Hamdi Alakkad.  In Section 4.4, item 2, the paper discusses finite-dimensional certificates for the diametral pressure functional.  The statement spans a page break in the PDF; the stitched screenshot below shows the relevant passage at full page width.

\begin{center}
\includegraphics[width=\textwidth]{figures/open_problem_crop.png}
\end{center}

The specific sentence tested here is the assertion that a finite family of unit vectors in a Hilbert space with sufficiently small mutual coherence gives finite-level certificates which force
\[
  \mathbf P(C_1,x_\infty)>0.
\]
This packet does not address the main reflexive fixed point property conjecture in the paper.  It only refutes this finite-certificate-to-global-pressure implication.

\section*{Definitions}

Following Definition 4.1 of the source paper, for a bounded set \(C_1\), a point \(x_\infty\), and \(\Delta=\operatorname{diam}(C_1)\), set
\[
  \Phi_k(C_1,x_\infty)
  =
  \sup_{y_1,\ldots,y_k\in C_1}
  \inf_{\|a\|_1=1}
  \left\|
  \sum_{i=1}^k a_i\,\frac{y_i-x_\infty}{\Delta}
  \right\|.
\]
The global pressure is
\[
  \mathbf P(C_1,x_\infty)=\inf_{k\ge 1}\Phi_k(C_1,x_\infty).
\]

\section*{Counterexample}

\begin{theorem}
Finite low-coherence Hilbert certificates through \(k\le m\) do not force positive global pressure.
\end{theorem}

\begin{proof}
Let \(H=\mathbb R^2\) with its Euclidean norm, let \(x_\infty=0\), and let
\[
  C_1=\{e_1,e_2\}.
\]
The vectors \(e_1,e_2\) are unit vectors and their mutual coherence is
\[
  \mu=|\langle e_1,e_2\rangle|=0<1=\frac{1}{2-1}.
\]
Also \(\Delta=\operatorname{diam}(C_1)=\|e_1-e_2\|=\sqrt 2\).

For \(k=2\), choosing the tuple \((e_1,e_2)\) gives
\[
  \inf_{|a_1|+|a_2|=1}
  \left\|
    \frac{a_1e_1+a_2e_2}{\sqrt2}
  \right\|
  =
  \frac{1}{\sqrt2}
  \inf_{|a_1|+|a_2|=1}(a_1^2+a_2^2)^{1/2}
  =
  \frac{1}{2}.
\]
Thus the advertised finite-level phenomenon is present: \(\Phi_2(C_1,0)\ge 1/2\), and of course \(\Phi_1(C_1,0)=1/\sqrt2\).

However, for \(k=3\), every 3-tuple chosen from the two-point set \(C_1\) repeats a point.  If \(y_p=y_q\), choose coefficients
\[
  a_p=\frac{1}{2},\qquad a_q=-\frac{1}{2},\qquad a_i=0\quad (i\ne p,q).
\]
Then \(\|a\|_1=1\) and
\[
  \sum_{i=1}^3 a_i\,\frac{y_i-x_\infty}{\Delta}=0.
\]
Hence the inner infimum is zero for every 3-tuple, so \(\Phi_3(C_1,0)=0\).  Consequently
\[
  \mathbf P(C_1,0)=\inf_{k\ge1}\Phi_k(C_1,0)=0.
\]
Therefore the finite low-coherence lower bound for \(k\le m\) does not force positive global pressure.
\end{proof}

\begin{remark}
The obstruction is not special to two points.  If \(C_1\) is finite with \(|C_1|=m\), then every \(k\)-tuple with \(k>m\) repeats a point, and the same cancellation gives \(\Phi_k(C_1,x_\infty)=0\).  This is also consistent with Lemma 2.5 of the source paper.  The error is the passage from finite-level certificates to a uniform lower bound over all \(k\).
\end{remark}

\section*{Verification code}

The packet includes a small exact verifier:
\[
\texttt{code/verify\_pressure.py}.
\]
It enumerates tuples from \(C_1=\{e_1,e_2\}\) and computes the exact values
\[
  \Phi_1=\frac{1}{\sqrt2},\qquad \Phi_2=\frac{1}{2},\qquad \Phi_k=0\quad(k\ge3).
\]
Run it with:
\begin{verbatim}
conda run --no-capture-output -n sandbox python \
  code/verify_pressure.py
\end{verbatim}

\section*{Novelty and limitation}

This is a counterexample to a stated implication in the paper's programmatic discussion.  It should not be counted as a solution to the author-labeled open problems in Section 4.6, and it does not disprove the main reflexive fixed point property conjecture.

\end{document}
